% !TeX root = ../papers/nc-p3.tex
\section{全量子背景の三つの余指標表と有限重み再帰}
\label{sec:quantum-universal-chambers}
前二節で除いた面重複、根1、可換点も含め、軌道閉包に量子点を持つ範囲全体で、
候補の生成・正則性・添字保持同型を扱う。以下の三表では任意の体要素を自由欄に入れ、
線型方程式の右辺だけを検査する。退化先の量子点は定理\ref{thm:quantum-indexed-all}の
三パラメータと有限置換で指定する。

\begin{theorem}\label{thm:quantum-universal-chambers}
任意の量子背景 $w$ を固定する。以下の手順で受理する窓が、軌道閉包に $w$ を持つ
全幾何軌道の代表を尽くす。全て標準Koszul左右AS正則である。文字の番号付けを全て許して、
必要な余指標と正の整数重みの上限は次の三行だけである。
\[
\begin{array}{c|c|c}
 \lambda & \text{正の多重次数の個数}&\text{最大重み}\\ \hline
 (0,1,4,6)&17&13\\
 (0,3,5,9)&17&19\\
 (0,6,8,9)&20&13
\end{array}
\]
表内・表間の同型は、下記の横断面上で一つの四次行列の安定化群作用により判定できる。
非零台や安定化群の一般性を仮定しない。
\end{theorem}
\begin{proof}[安定化群と横断面]
同型と正則性に必要な基底は射影的に取る。まず全行積が1となる量子表示を使う。
さらに二文字 $a,b$ を含む三角積が全て1なら、その文字間の交換係数を1とし、他文字との交換係数も
二文字で一致するようにする。実際、定理\ref{thm:quantum-indexed-all}の証明によれば、
面値1が二つなら一つの二文字組、四つなら可換表示にできる。二文字を揃えた後、行積を
$P_j$ として対角twist倍率の四乗を $P_j^{-1}$ に取る。二文字組では同じ根を選び、
倍率の全積は組の外の一文字の四乗根の選択で1にできる。面値1が高々一つでも同じ行積調整を使う。
これで四位置に同じ対角トーラス $T$ が作用し、背景の量子体積は一位置回転で符号を変える。

四基底群を $G=(\mathrm{PGL}_4)^4$、射影安定化群を $H$ とする。
点対応の最初の射影変換で $H$ は忠実に表される。面値1が0または1個なら座標点とその右対応が内在的である。
2個なら二文字の全射影直線上でも右対応が同じ点であり、
外の二座標点とそれらから各内側座標点への直線上での右対応が残る対角倍率の比を決める。
4個なら右対応は全射影空間で恒等である。従って位置0の行列は、
同一交換係数の組の内部での可逆ブロック、各一文字の倍率、三角積を保つ置換の積で尽きる。
これらは全て持ち上がる。ブロックは四位置で共通であり、置換には量子同型判定の位置別倍率を付ける。
位置0を固定した持上げは他位置でも射影的に一意である。
$T$ は $H^\circ$ の極大トーラスである。

射影基点直線の $H$ 指標を除いてテンソル表現を正規化する
（線型表現は原空間と基点直線の双対との積に誘導される）。
$H$ は簡約的なので $H$ 不変線型形式 $f(w)=1$ と、
$f=1$ の中で $G$ の射影軌道接に横断する $H$ 不変補空間 $N$ を取れる。
巡回も同時に扱うときは一位置回転を加えた有限拡大の安定化群で選ぶ。
これらの選択は有限線型代数で可能である。例えば $f$ と軌道接への恒等な射影について、
既知ブロック群のLie環と有限置換、必要なら回転との可換性を線型条件として解く。
軌道接は四位置で各行列単位を $w$ に作用させ、射影方向へ写せば求まる。

二隣接階数が6の局所閉スキームを $Z$ とする。量子点のpolystabilityは
定理\ref{thm:quantum-face-resonance}の台の重心の証明が全 $q$ に適用できる。
Lunaの定理を $w$ で非零の不変斉次式のアフィン開集合に適用する。
$S=(w+N)\cap Z$ の基点ファイバーから $G$ 側の全基点ファイバーへ、
$H$ 軌道と $G$ 軌道の全単射を得る。詳しくは、周囲の射影空間における
$G\mathbin{\times}^H(w+N)$ の局所写像が基点でエタールであり、$Z$ で基底変換して切片とする。
強エタールな飽和開集合に縮めると商との正方形がファイバー積であり、閉点の商ファイバー全体を含む。
従ってこの主張は形式的な弧だけに限定されない。
逆に横断面で $w$ に流れる全点は、使った各 $H$ 不変開集合の中に入る。
Hilbert--Mumford判定と極大トーラスの共役性から、$H$ の作用後に
$n\in N$ が $T$ のある余指標で正の成分のみを持つことが必要十分である。
\end{proof}

\begin{proof}[三表の全被覆]
語 $v$ の指標は $\beta(v)=(\#_0v,\ldots,\#_3v)-(1,1,1,1)$ である。
非零な34通りの超平面を、定数ベクトルを除いた余指標の三空間で使う。
壁上の正台は隣の室の正台に含まれるため包含極大な正台のみを残す。
列挙は整数演算で行える。各法線の第四座標を除き、独立な順序付き二法線 $a,b$ ごとに
$r=a\times b,\ e=a\times r$ と置く。各符号 $i,j,k\in\{\pm1\}$ に対して、
別の法線 $v$ と $(ir,je,ka)$ との三内積のうち最初の非零値が正かを記録する。
これは一つの半直線、その上の二次元面、その外側の室を通る微小変位であり全室を尽くす。
実際、法線が三空間を張るので各全次元室の閉錐には稜があり、その稜を含む面を選べる。
重複を除くと336室、包含極大は17文字の48個と20文字の24個となる。
\path{computations/verify_quantum_weight_recursion.py} はその72個が表の三余指標の全置換に
一致することを確認する。余指標を極大台の代表へ換えても元の全非零成分は正のままである。
\end{proof}

\begin{proof}[重み順の復元と判定]
選んだ室の正の $N$ の座標だけを使う。各切断 $M$ の行列を順序対で表示する。
行・列とも $i<j$ の六対 $P$ を主部分に取る。量子背景では各行は相補対の列にちょうど一つの
非零体積係数を持ち、$M_{PP}$ は可逆である。この行列式の指標は零である。
実際、六対の文字総数は各添字3で、行と列を足して六つの体積指標を引くと零となる。
正の変位だけを許すと行列式は従って背景値のままである。残る十対を $E$ と置く。
求める方程式は各切断の百成分
\[
 B(M)=M_{EE}-M_{EP}M_{PP}^{-1}M_{PE}=0
\]
である。分母が非零定数なのでこの表の上では多項式となる。
行 $u$、列 $v$ の成分の指標は語 $uv$ の指標である。
これは六次主小行列を縁取る七次行列式をその主行列式で割る式からも分かる。

二切断を合わせた $B$ の基点での微分を $L$ とする。各多重次数 $\beta$ ごとに
$N_\beta=K_\beta\oplus D_\beta$, $K_\beta=\ker L_\beta$ と線型分解する。
整数重み $d=1,\ldots,\max\langle\beta,\lambda\rangle$ を順に処理する。
各 $\beta$ の自由欄に $z_\beta\in K_\beta$ を指定し、既に復元した
$n_{<d}$ を $B(w+n_{<d})$ に代入する。その $\beta$ 成分の負を $b_\beta$ とする。
$b_\beta\notin\operatorname{im} L_\beta$ なら棄却し、そうでなければ
\[
 n_\beta=z_\beta+(L_\beta|_{D_\beta})^{-1}b_\beta
\]
を加える。$N_\beta=0$ のときにも右辺の整合性を検査する。
多項式の非線型項に重み $d$ 以上の入力があれば、もう一つの正の入力を要するため
出力の重みは $d$ より大きい。従ってこの処理は重み $d$ の全式を厳密に解く。
非正重みの式は背景の値であり、最後の重みの後に方程式は残らない。
帰納的に、受理する各欄の値がただ一つのテンソル $w+n$ を与え、表内の全解がこの形となる。
両階数は常に6で、背景への一径数退化に命題\ref{prop:unrestricted-good-reduction}を適用できる。

異なる背景は定理\ref{thm:quantum-indexed-all}によって先に比較する。背景を同じ表示に揃えた後、
二つの復元値が同型となる条件は $n'=hn$ ($h\in H$) である。
有限置換を試し、残る四次のブロック行列に逆行列式用の一変数を付けて
多項式イデアルが単位イデアルかを消去すれば、非可換な安定化群でも判定が終わる。
周期は同じ検査で一位置・二位置の巡回を比較して決まる。
同じ商ファイバー内での軌道閉包も切片に移せる。ブロック群の各成分で
$y=hn$ のグラフから群変数を消去して得るイデアルを $n'$ に代入すればよい。
座標環の消去定理によりこの操作が像のZariski閉包を与える。トーラスだけの場合は
既出の整数核と台の面の検査に簡約する。

検算では三面一致の有理量子点で九方向のうち同じ余指標に正の全方向を同時に復元した。
可換点では、接空間の自由欄で実際に棄却する入力と、非線型な強制補正を一回
要して受理する入力の両方を有理数上で照合した。
\end{proof}
